\section{Thảo luận}\label{sec:discussion}

\paragraph{Điều tách quyết định khỏi diễn đạt mua được, và điều nó không mua được.}
Kết quả nhất quán ở ba thực nghiệm độc lập về cùng một điều: những tính chất mà bài báo gọi là bất biến (không lộ đáp án, không khen sai, không dạy trước, không để mô hình quyết định đúng--sai) giữ được \emph{bằng cấu trúc} chứ không bằng lời nhắc. Chatbot dùng lời nhắc thường vi phạm chúng ở một nửa số lượt thử thách; một lời nhắc soạn kỹ làm hiếm đi nhưng không loại hẳn, và lần vi phạm còn lại là loại nguy hiểm nhất, lại là loại một giám khảo cùng kiểu không thấy. Nhưng cấu trúc không mua được sự tự nhiên. Trên thang của giám khảo, B2 nghe như một cô giáo hơn CAVAT, và các lượt kém nhất của CAVAT là những lượt mà một chính sách tất định xử lý vụng: nói lại đúng một câu khi thang cạn, xác nhận không hỏi vì sao vì băm của mã câu rơi ngoài $30\%$, đưa một phản ví dụ gắn với lỗi chứ không gắn với câu. Đánh đổi này không phải ngẫu nhiên: mỗi bảo đảm là một điều mô hình không được tự do làm.

\paragraph{Ràng buộc không miễn phí, và đó là một kết quả đáng giá hơn một con số tốt.}
Trực giác thường là ràng buộc theo đồ thị tri thức làm truy xuất ``thông minh hơn''. Dữ liệu của chúng tôi nói khác: phạm vi chặt khi chưa có định tuyến cho Hit@3 trên tập giữ riêng chỉ $0{,}45$, so với $0{,}90$ của RAG phẳng, vì nó đẩy ra ngoài chính các đoạn trả lời câu hỏi về bài cũ; định tuyến khôi phục lại, và kết quả cuối ngang một RAG phẳng. Lợi ích có thật nhưng thuộc loại khác: không có đoạn bài sau nào lọt vào câu trả lời (bất biến B5), điều mà RAG phẳng làm ở khoảng một phần ba số câu. Một người xây hệ sau này nên kỳ vọng ràng buộc phạm vi là một khoản chi để mua an toàn, phải được bù bằng một cơ chế như định tuyến, và nên đo nó như vậy.

\paragraph{Bộ kiểm chứng là thành phần chính.}
Qwen viết câu bị chặn ở nhiều chỗ, và mỗi mã lỗi trong Bảng~\ref{tab:validator} có một ví dụ thật: thuật ngữ ngoài đoạn nguồn, một câu nối nói ý sai của trẻ như thể là ý của gia sư, số không có trong đề, và $33\%$ gợi ý do mô hình soạn không qua kiểm (Bảng~\ref{tab:drafts}). Nếu coi mô hình là thành phần chính và bộ kiểm chứng là bảo hiểm, hình dung đó ngược với quan sát: đa số lời nháp của mô hình, ở các tác vụ có ràng buộc, cần bị chỉnh hoặc bỏ. Điều này cũng cho thấy một chi phí ẩn của hướng thiết kế này: bộ kiểm chứng và các câu mẫu là nội dung phải được bảo trì cùng với chương trình.

\paragraph{Giám khảo mô hình ngôn ngữ cần một bộ đo tất định đi kèm.}
Ca ``đã viết đúng số thập phân rồi'' là một ví dụ cụ thể và nhỏ của một điều đã được nêu chung~\cite{zheng2023}: giám khảo cùng loại với hệ được chấm có xu hướng chấp nhận chính loại lỗi của nó. Không thể rút ra từ một lượt rằng giám khảo kém; có thể rút ra rằng với các thuộc tính mà một luật kiểm được (lộ đáp án, một câu hỏi, số), nên đo bằng luật, và dùng giám khảo cho những thứ chỉ giám khảo mới thấy (giọng điệu, tính hợp tuổi).

\paragraph{Diễn giải khai báo làm cho đánh giá kiểm toán được.}
Trong mô hình người học, mọi chỗ khung không nói được đánh dấu DI và có số hiệu; nhờ đó phân tích độ nhạy (Bảng~\ref{tab:sensitivity}) là một phép toán trên một danh sách có sẵn, và nó chỉ ra đúng nơi nên đầu tư hiệu chỉnh: ngưỡng giải thích. Một DI sai cũng được sửa có ghi lại: DI-55 bản đầu trừ điểm lời giải thích chung chung và làm kẹt học sinh mẫu ở tiết đầu, và thực nghiệm đã đổi nó. Một hệ thống chỉ dựa vào ``mô hình chấm'' không có chỗ để bám như vậy.

\paragraph{Nội dung là điểm yếu thật.}
Hai điểm yếu rõ nhất trong dữ liệu (lời gợi ý lặp lại nguyên văn ở $4\%$ lượt; phản ví dụ không bám câu) đều không phải lỗi của mô hình mà của \emph{nội dung đã duyệt}: thang gợi ý không đủ bậc cho mục tiêu mà luật chọn, phản ví dụ gắn với lỗi chứ không gắn với câu. Chúng sửa được bằng nội dung, và đó cũng là điểm mà người giáo viên (chứ không phải kỹ sư) có thể đóng góp.

\paragraph{Điều chưa biết.}
Chúng tôi chưa biết liệu các bảo đảm trên làm học sinh \emph{học tốt hơn}. Có thể một gia sư tự nhiên hơn nhưng đôi khi lộ đáp án vẫn cho kết quả học tập ngang hoặc tốt hơn, nếu sự tự nhiên giữ trẻ ở lại với bài; ngược lại cũng có thể một gia sư cứng hơn gây nản. Bastani và cộng sự~\cite{bastani2025} cho thấy các ràng buộc kiểu gia sư có thể giảm tác hại của truy cập không hạn chế; chúng tôi không có bằng chứng ở chiều kết quả cho hệ này. Hướng tự nhiên nhất để tách hai câu hỏi là một thực nghiệm trong đó phần diễn đạt của B2 được dùng \emph{bên trong} rào chắn của CAVAT (chính sách, bộ chấm, bộ kiểm chứng), để xem có giữ được giọng điệu mà không mất bảo đảm hay không.
